\documentclass[10pt,a4paper]{article}
\usepackage[margin=0.78in]{geometry}
\usepackage{amsmath,amssymb,booktabs,array}
\usepackage[hidelinks]{hyperref}
\setlength{\parindent}{0pt}
\setlength{\parskip}{5pt}
\newcommand{\head}[1]{\par\medskip\textbf{#1}\par}
\newcommand{\Try}{\mathcal L}
\begin{document}
\begin{center}
{\large\bfseries A conservative allocation completion for the VSL lane candidate}\\[3pt]
2026-10-21 --- proposed adaptation; no model adoption or traffic experiment
\end{center}

\textbf{Result and boundary.} A small explicit allocator can satisfy common
traffic and receiving budgets on the three-lane geometry. The choice between
unsuccessful lane-change attempts continuing forward and committed directional
fractions changes queue behavior. This note specifies the former and compares
the latter. It establishes conditional algebraic admissibility, not a complete
VSL/capacity-drop/task qualification or a calibrated lane-change law.

\head{1. State, geometry and admissible inputs}
Arrange lanes as $E,A,B$: ending lane, adjacent continuing lane, far continuing
lane. Physical cell row $i$ sends to row $i+1$ through a storage-free allocation
node. An allowed movement goes straight or to an adjacent lane; it advances
exactly one row. All boundary flows are evaluated from the old time level
before any stock is updated. This diagonal transport is motivated by Carey
et al.~\cite{carey}; it replaces the principal source's same-cell lateral step
followed by adjusted longitudinal flows~\cite{subraveti}.

For one boundary, let $a$ index donor lanes and $j$ receiver lanes. Use amounts
per fixed time step, not rates:
\begin{equation}
 d_a=\Delta t\,D_{i,a},\qquad r_j=\Delta t\,R_{i+1,j},\qquad
 0\leq d_a\leq n_{i,a},\quad
 0\leq r_j\leq J_{i+1,j}-n_{i+1,j}. \label{eq:inputs}
\end{equation}
$n$ is continuous vehicle-equivalent stock, $J$ jam storage, and $D,R$ have
units of vehicle-equivalents per unit time. Thus $d,r,n,J$ have the same mass
unit. Nothing is rounded. Source sending, receiving and any capacity-drop
branches must actually satisfy \eqref{eq:inputs}; this is a condition, not a
result of the allocator. Standard fixed-grid CTM bounds include
\begin{equation}
 v_{\max}\Delta t/\Delta x\leq1,\qquad
 w_{\max}\Delta t/\Delta x\leq1. \label{eq:cfl}
\end{equation}
For a triangular FD, $D\leq v_{\max}k$ and
$R\leq w_{\max}(k^{\rm jam}-k)$ then imply the stock bounds. Modified source
branches and the eventual Han-type sending restriction~\cite{han} need their
own compatibility check; \eqref{eq:cfl} alone does not certify every branch.

\head{2. Proposed translation of incentives into requests}
Let $P_{aj}\geq0$ be the principal source's finite lateral preference for an
allowed adjacent direction, before its disputed realization equations.
Mask forbidden directions first. For a continuing donor define
\begin{align}
 s_a&=\sum_{j\in\Try_a}P_{aj}, &
 b_{aj}&=\frac{P_{aj}}{\max(1,s_a)}\quad(j\in\Try_a),\nonumber\\
 b_{aa}&=1-\sum_{j\in\Try_a}b_{aj}, & m_{aj}&=d_a b_{aj}. \label{eq:requests}
\end{align}
$\Try_a$ is the set of allowed adjacent targets; $m_{aa}=d_ab_{aa}$.
Hence $m_{aj}\geq0$ and $\sum_jm_{aj}=d_a$. Existing valid fractions are
unchanged; only a jointly oversubscribed donor is rescaled. This normalization
is a proposed interface rule, not a native source equation. For an empty
donor all requests are zero; no empty-density preference ratio is evaluated.

At the permanent end there is no $E$ receiver: remove movements into it.
Set $b_{EA}=1$ for the last E donor. No $E\to B$ movement is permitted.
Before the end, $E\to E$ remains a legal forward fallback even if its initial
request happens to be zero. Fractions are evaluated anew from the next state;
no persistent route/intention records are introduced.

\newpage
\head{3. Concrete candidate: allocate attempts, then allow forward fallback}
\textbf{First pass: share each receiving budget.}
For each target $j$, fix positive priorities $p_{aj}$ for its nonzero request
contributors. The proposal uses the donor's \emph{reference geometric lane
capacity} (times $\Delta t$, if a mass scale is desired), independently of
advice and current demand. Equal reference capacities give equal priorities.
This uses the capped fixed-priority sharing principle of Carey R3 and the
MISO node literature~\cite{carey,wright}; the three-lane composition is ours.
Define
\begin{equation}
 z_{aj}=\begin{cases}
 m_{aj},& \sum_h m_{hj}\leq r_j,\\
 \min(m_{aj},\lambda_jp_{aj}),&\sum_h m_{hj}>r_j,
 \end{cases}
 \quad\text{with}\quad
 \sum_h\min(m_{hj},\lambda_jp_{hj})=r_j. \label{eq:share}
\end{equation}
No demand means zero flow. At zero supply take $\lambda_j=0$. In the
constrained positive case the flow vector is unique: increase the common
allocation level until all the available receiving budget is used. A request
that is fully served drops out of further sharing, leaving space for others.
In particular, $0\leq z_{aj}\leq m_{aj}$ and
$\sum_a z_{aj}\leq r_j$. A small or absent request cannot retain an unused
reserved entitlement.

\textbf{Second pass: use otherwise unused forward space.}
Define residual donor traffic and receiver space after \eqref{eq:share}:
\begin{equation}
 e_a=d_a-\sum_jz_{aj}\geq0,\qquad
 h_j=r_j-\sum_az_{aj}\geq0. \label{eq:residuals}
\end{equation}
If the same-lane connection exists, send
\begin{equation}
 y_a=\min(e_a,h_a),\qquad
 f_{aj}=z_{aj}+\mathbf1_{\{j=a\}}y_a. \label{eq:fallback}
\end{equation}
When it does not exist, set $y_a=0$. Each receiver has at most one such
same-lane fallback donor. There is no second lateral reallocation and no
same-step movement of newly arriving stock. The two passes decide one set
of flows; they are not two physical time steps.

The elementary proposal gives one entering mass unit one receiving-budget
unit. It sets Carey's extra headway multiplier to $\alpha=1$ rather than
adding another merging-loss law. This is an explicit modeling simplification.
Extra gap consumption and slow-entry recovery are not thereby reproduced;
the retained congestion/capacity-drop mechanism still needs qualification.

\head{4. Strong rival: keep directional fractions as commitments}
Instead feed $d_a$, $b_{aj}$, $r_j$ and fixed input priorities into the
published general junction allocation framework~\cite{wright}. With full
FIFO within an input lane, its directions are restricted together:
\begin{equation}
 f_{aj}=b_{aj}F_a,\qquad 0\leq F_a\leq d_a,
 \qquad \sum_a b_{aj}F_a\leq r_j. \label{eq:fifo}
\end{equation}
The source's priority/redistribution rule, not these constraints or maximum
flow alone, completes the multi-input solution. This is the alternative's
published algorithmic contract; it is not an arbitrary LP tie break.
Blocked turns can hold back all directions from their input lane. No-FIFO
instead caps each movement by its own $m_{aj}$; this still does not add the
forward fallback of \eqref{eq:fallback}. These fractions need only be fixed
within one evaluation, not for a driver's entire trip.

Both candidates require explicit normalized incentive and boundary mappings.
Neither can silently import its source's general-node theorems into the full
VSL plant. In particular the preferred candidate is \emph{not} the Wright
MIMO solver. Demand-proportional priorities have a source invariance warning
\cite[Remark 3.1]{wright}; using fixed priorities avoids that particular
choice, but is not a proof of the preferred composition's nodal invariance.

\newpage
\head{5. Conditional admissibility proof for the preferred candidate}
From \eqref{eq:requests}--\eqref{eq:fallback},
\begin{align}
 f_{aj}&\geq0, &
 \sum_j f_{aj}&=\sum_jz_{aj}+y_a\leq d_a,\nonumber\\
 \sum_a f_{aj}&=\sum_az_{aj}+y_j\leq r_j. \label{eq:budgetproof}
\end{align}
For a receiver without a legal same-lane donor, its added $y_j$ is zero.
Thus every donor and receiver is shared once even with two lateral directions
and simultaneous opposing changes. Individual desired fractions are not
independent permissions to reuse mass. In the absence of lateral requests,
the first pass gives the ordinary lane CTM flow $\min(d_a,r_a)$ and fallback
adds zero.

Let $O_{i,a}$ sum all movements from cell $(i,a)$ and $I_{i,a}$ sum all
movements entering it from row $i-1$. The stock update is
\begin{equation}
 n^{+}_{i,a}=n_{i,a}-O_{i,a}+I_{i,a}. \label{eq:update}
\end{equation}
Because $O_{i,a}\leq d_{i,a}\leq n_{i,a}$ and
$I_{i,a}\leq r_{i,a}\leq J_{i,a}-n_{i,a}$,
\begin{equation}
 0\leq n_{i,a}-d_{i,a}\leq n^{+}_{i,a}
 \leq n_{i,a}+r_{i,a}\leq J_{i,a}. \label{eq:storage}
\end{equation}
Every internal movement appears once as a departure and once as an arrival,
so summing \eqref{eq:update} conserves mass up to recorded external flows.
Only allowed connections carry requests or fallback. Unserved final-E mass
remains in the stock update; reaching the end is not an exit.

\textbf{What replaces the earlier obstruction?} There is no adjusted
receiving term $R+L_{\rm in}-L_{\rm out}$ and no negative residual demand.
Flows use nonnegative budgets, and realized outflow cannot exceed their
shared donor. A departure's released storage becomes available through the
next old-state evaluation. This changes the time/space stencil rather than
asserting a sign correction to the native formula. The earlier source
counterexamples remain unchanged; they are not relabeled as trajectories
of this new model.

\head{6. Local sensitivity and important missing guarantees}
For a constrained receiver, fix the active sets $H$ of fully served requests
and $U$ of requests limited by their priority entitlement. Then
\begin{equation}
 \lambda_j=
 \frac{r_j-\sum_{a\in H}m_{aj}}{\sum_{a\in U}p_{aj}},\qquad
 z_{aj}=\begin{cases}m_{aj},&a\in H,\\
 \lambda_jp_{aj},&a\in U.
 \end{cases} \label{eq:active}
\end{equation}
The denominator is positive. Hence the allocator is smooth on fixed positive
request active sets wherever its input maps are smooth. Minima, the donor
normalization threshold, empty requests and changes in these sets require
branch-aware derivatives. This is a local algebraic observation, not finite
difference validation, a mode-event treatment or useful-gradient evidence.

The source incentive and capacity-drop branches may themselves have
nonsmooth or discontinuous transitions; no complete derivative statement is
made about that joined model. Nor are continuum convergence, invariance or
global throughput maximization claimed for this two-pass adaptation. The
next page gives a concrete throughput limit. Admissibility alone does not
establish a realistic merge service or a useful VSL task.

\newpage
\head{7. Hand cases: direct answers rather than a rollout}
All numbers below are nonnegative allocator input amounts per step. They
test the proposed budget rules, not the existence of a complete density,
incentive and control trajectory with exactly these values.

\textbf{Two neighbors overbook one donor.} Set $d=10$ and two lateral
preferences $P_1=P_2=7/10$. Equation~\eqref{eq:requests} gives
\begin{equation}
 (m_1,m_2,m_{\rm own})=(5,5,0).
\end{equation}
With no other contenders and receiving budgets $(2,8,10)$, the first pass
serves $(2,5,0)$. The remaining three use the own-lane connection:
\begin{equation}
 (f_1,f_2,f_{\rm own})=(2,5,3),\qquad \sum_jf_j=10.
\end{equation}
No negative residual, mass multiplication or unused donor quota is hidden.

\textbf{Failed attempt versus committed movement.} Before the lane end let
$d_E=10$, $b_{EA}=b_{EE}=1/2$, $r_A=0$, $r_E=10$, with no other contenders.
Then
\begin{center}
\begin{tabular}{lrr}
\toprule
Rule & $E\to A$ & $E\to E$\\
\midrule
Attempts, then forward fallback & 0 & 10\\
Fixed fractions, full FIFO & 0 & 0\\
Fixed fractions, no FIFO, no fallback & 0 & 5\\
\bottomrule
\end{tabular}
\end{center}
At the permanent end remove $E\to E$ and set $b_{EA}=1$: with $r_A=0$
all ten remain in the last E cell for either main candidate.

\textbf{Cooperation under constrained receiving.} Let E and A each request
10 units into A, with equal positive reference priorities and $r_A=8$.
Equation~\eqref{eq:share} serves $(4,4)$. If E requests only 1 and A still
requests 10, it serves $(1,7)$: unused entitlement is returned. This is fixed
priority with demand caps, not a fixed realized merging fraction.

\textbf{A real limitation: the preferred rule need not maximize total flow.}
Let $d_E=d_A=10$, $b_{EA}=4/5$, $b_{EE}=1/5$, and $b_{AA}=1$ before the
lane end. Set $r_A=r_E=10$ and use equal reference priorities. The first
pass serves $E\to A=5$, $A\to A=5$, $E\to E=2$; E's fallback adds 3.
The final total is $15$. If the eight lateral requests could instead all
be waived before sharing, sending both complete donors forward would give
$20$. Our rule gives initial requests precedence; the later fallback does
not revoke space they already claimed. A maximal total-flow objective with
freely adjustable desires would be a different behavioral model.

\textbf{Decision implication.} The preferred rule offers a complete,
deterministic allocation component with transparent limitations. Accepting
its opportunity meaning and priority sequence is a substantive modeling
choice. This proof does not choose it on the user's behalf or close the
remaining task qualification.

\begin{thebibliography}{9}\setlength{\itemsep}{0pt}
\bibitem{carey} M. Carey, C. Balijepalli, and D. Watling, ``Extending the cell
transmission model to multiple lanes and lane-changing,'' \emph{Netw. Spat.
Econ.}, vol. 15, pp. 507--535, 2015, doi: 10.1007/s11067-013-9193-7.
Methods read in the institutional 2013 author manuscript, pp. 4--12.
\bibitem{wright} M. A. Wright, G. Gomes, R. Horowitz, and A. A. Kurzhanskiy,
``On node models for high-dimensional road networks,'' \emph{Transp. Res. B},
vol. 105, pp. 212--234, 2017, doi: 10.1016/j.trb.2017.09.001.
Methods read in arXiv:1601.01054 HTML, secs. 2.1 and 3.1--3.3.
\bibitem{subraveti} H. H. S. Nagalur Subraveti, V. L. Knoop, and B. van Arem,
``First order multi-lane traffic flow model -- an incentive based macroscopic
model to represent lane change dynamics,'' \emph{Transportmetrica B},
vol. 7, no. 1, pp. 1758--1779, doi: 10.1080/21680566.2019.1700846.
\bibitem{han} Y. Han \emph{et al.}, ``Resolving freeway jam waves by discrete
first-order model-based predictive control of variable speed limits,''
\emph{Transp. Res. C}, vol. 77, pp. 405--420, 2017,
doi: 10.1016/j.trc.2017.02.009.
\end{thebibliography}
\end{document}
